\section{Agriculture de précision :  concepts fondamentaux}

\subsection{Définition et principes de l'agriculture de précision}
L'agriculture de précision est une stratégie de gestion agricole fondée sur la prise en compte de la variabilité spatiale et temporelle des parcelles afin d'améliorer la durabilité agricole \cite{monteiro_precision_2021}. Contrairement aux méthodes conventionnelles, qui appliquent des interventions uniformes sur l'ensemble d'une exploitation, elle adapte les apports en eau, en fertilisants et en amendements aux besoins réels des cultures, en fonction des caractéristiques du sol, des conditions climatiques et du stade de développement des plantes.

Cette méthode repose sur la collecte, le traitement et l'analyse de données issues de différentes technologies, notamment les capteurs de sol, les stations météorologiques, les systèmes de géolocalisation (GPS), la télédétection, les drones et les analyses agronomiques \cite{shaheb_precision_2022}. Les informations recueillies permettent d'évaluer l'état des cultures et de leur environnement, facilitant ainsi une prise de décision plus précise.

L'agriculture de précision s'articule autour de trois principes fondamentaux : observer, analyser et agir \cite{noauthor_precision_2026}. Il s'agit d'abord de décrire la variabilité des parcelles, puis analyser les données afin d'identifier les besoins des cultures et enfin, d'ajuster les interventions (l'irrigation et la fertilisation), en fonction des conditions réelles du champ. 
L'objectif est d’améliorer l'efficacité de la gestion des ressources, de réduire les impacts environnementaux et de maintenir, voire augmenter les rendements.

\subsection{Origines et évolution du concept}
Le concept d’agriculture de précision s’est développé progressivement pour mieux prendre en compte la variabilité des sols, des cultures et des conditions de  l'environnement sur une même parcelle. 
Initialement fondée sur des observations et des mesures des propriétés du sol et des rendements, cette méthode a évolué avec l’introduction des outils de cartographie, du GPS et des systèmes d’information géographique, permettant une gestion différenciée des parcelles.

Les progrès des capteurs, de la télédétection et de l’informatique ont ensuite facilité la collecte et l’analyse des données météorologiques et agronomiques. 
Aujourd’hui, l’intégration de technologies connectées, notamment les \textbf{réseaux de capteurs et l’Internet des objets (IoT)} favorise un suivi plus régulier des parcelles et un ajustement des interventions en fonction des besoins réels des cultures. 
Cette évolution vise à améliorer l’utilisation des ressources, à \textbf{réduire les pertes et à préserver la productivité agricole}.

\subsection{Agriculture de précision versus agriculture traditionnelle}
L'agriculture traditionnelle repose généralement sur la gestion uniforme des parcelles, avec des apports en eau et en intrants 
déterminés selon l'expérience de l'agriculteur, un calendrier fixe ou l'observation visuelle des cultures. 
En revanche, l'AP s'appuie sur des capteurs et des données météorologiques afin d'adapter les interventions face aux variations du sol, du climat et des besoins réels des cultures.

Elle vise à améliorer l'utilisation de l'eau et des intrants, à limiter les pertes et à maintenir, voire augmenter les rendements. 
Toutefois, son adoption nécessite des compétences techniques, des équipements et un investissement initial important. Elle cherche donc à améliorer les pratiques traditionnelles en les rendants plus éfficaces.

\subsection{Enjeux et intérêts pour l'agroécologie}
L'agriculture de précision contribue à l'agroécologie en optimisant l'utilisation des intrants, des amendements organiques et de l'eau selon les caractéristiques du sol et les besoins réels des cultures. 
Elle permet ainsi de réduire le gaspillage des ressources et les impacts environnementaux, tout en améliorant l'efficacité des pratiques agricoles. 
Dans les zones confrontées à la raréfaction de l'eau et à la dégradation des sols comme le Sine Saloum, elle constitue une méthode pertinente pour concilier productivité agricole, préservation des ressources naturelles et durabilité des systèmes de production.
